\VerdierPhysicalPageAnchor{0044}%
\VerdierSemanticAnchor{ch1-1-4}
\subsection{Changement de groupes}

\VerdierSemanticAnchor{ch1-1-4-1}
\noindent\textbf{1.4.1.}\quad
Soit $f\colon H\to G$ un homomorphisme de groupes. On associe à toute
$G$-catégorie $(\mathcal C,G,T,c)$ une $H$-catégorie
$(\mathcal C,H,T_f,c_f)$ en faisant opérer $H$ sur $\mathcal C$ par
l'intermédiaire de $f$, \emph{i.e.} en posant :

\VerdierSemanticAnchor{ch1-1-4-1-eq1}
\begin{equation}
  \begin{aligned}
    T_f(h)&=T\bigl(f(h)\bigr),
      & h&\in H,\\
    c_f(h,h')&=c\bigl(f(h),f(h')\bigr),
      & h,h'&\in H.
  \end{aligned}
  \tag{1.4.1.1}
\end{equation}

\VerdierSemanticAnchor{ch1-1-4-1-gfunctor}
Étant donné de même un $G$-foncteur :
\[
  (F,m)\colon(\mathcal C,G,T,c)
  \longrightarrow(\mathcal C',G,T',c'),
\]
on lui associe un $H$-foncteur $(F,m_f)$ entre les $H$-catégories
correspondantes en posant :

\VerdierSemanticAnchor{ch1-1-4-1-eq2}
\begin{equation}
  m_f(h)=m\bigl(f(h)\bigr).
  \tag{1.4.1.2}
\end{equation}

\VerdierSemanticAnchor{ch1-1-4-1-morphisms}
Enfin, un morphisme de $G$-foncteurs compatible avec les opérations de
$G$ est compatible avec les opérateurs de $H$.

\VerdierSemanticAnchor{ch1-1-4-1-eq3}
Ce qui précède définit visiblement un foncteur, qu'on appelle le
\emph{foncteur de changement de groupes} :
\begin{equation}
  f^*\colon G\text{-}\operatorname{Cat}
  \longrightarrow H\text{-}\operatorname{Cat}.
  \tag{1.4.1.3}
\end{equation}

\VerdierSemanticAnchor{ch1-1-4-1-eq4}
Plus précisément, soient $\mathcal C$ et $\mathcal C'$ deux
$G$-catégories ; le foncteur $f^*$ définit une application :
\begin{equation}
  \begin{aligned}
    f^*(\mathcal C,\mathcal C')\colon\quad
      &\operatorname{Hom}_{G\text{-}\operatorname{Cat}}
        (\mathcal C,\mathcal C')\\
      &\longrightarrow
        \operatorname{Hom}_{H\text{-}\operatorname{Cat}}
        \bigl(f^*(\mathcal C),f^*(\mathcal C')\bigr).
  \end{aligned}
  \tag{1.4.1.4}
\end{equation}

\VerdierSemanticAnchor{ch1-1-4-1-homcategories}
Mais les ensembles
$\operatorname{Hom}_{G\text{-}\operatorname{Cat}}
(\mathcal C,\mathcal C')$ et
$\operatorname{Hom}_{H\text{-}\operatorname{Cat}}
\bigl(f^*(\mathcal C),f^*(\mathcal C')\bigr)$ sont des ensembles
d'objets de catégories
$\mathcal H\!om_{G\text{-}\operatorname{Cat}}
(\mathcal C,\mathcal C')$ et
$\mathcal H\!om_{H\text{-}\operatorname{Cat}}
\bigl(f^*(\mathcal C),f^*(\mathcal C')\bigr)$ dont les morphismes sont
les morphismes de foncteurs, compatibles avec les opérations des
groupes, et l'application $f^*(\mathcal C,\mathcal C')$ s'étend
naturellement en un foncteur :

\VerdierSemanticAnchor{ch1-1-4-1-eq5}
\begin{equation}
  \begin{aligned}
    f^*(\mathcal C,\mathcal C')\colon\quad
      &\mathcal H\!om_{G\text{-}\operatorname{Cat}}
        (\mathcal C,\mathcal C')\\
      &\longrightarrow
        \mathcal H\!om_{H\text{-}\operatorname{Cat}}
        \bigl(f^*(\mathcal C),f^*(\mathcal C')\bigr).
  \end{aligned}
  \tag{1.4.1.5}
\end{equation}

\VerdierSemanticAnchor{ch1-1-4-1-faithful}
Ce foncteur est fidèle.
